%---------------------------------------------------------------
% Formelverzeichnis mit Definition der Abkürzungen
%---------------------------------------------------------------
\header{Formelzeichen}

% Ein paar Vereinfachungen für Abkürzungen und so weiter
\newcommand{\zB}{z.\,B.} % sieht mit kleinem Abstand einfach besser aus
\newcommand{\uA}{u.\,A.} % Anwendung dann mit \zB\ etc.
\newcommand{\uU}{u.\,U.}
\newcommand{\dH}{d.\,h.} % \dh gibts leider schon, deshalb \dH

% Eigennamen als SmallCaps
\newcommand{\Matlab}{\textsc{Matlab}}
\newcommand{\Jacobi}{\textsc{Jacobi}}
\newcommand{\Kalman}{\textsc{Kalman}}
\newcommand{\Riccati}{\textsc{Riccati}}
\newcommand{\Pearl}{\textsc{Pearl}}

\newcommand{\abb}[1]{Abb.~\ref{fig:#1}}
\newcommand{\tab}[1]{Tab.~\ref{tab:#1}}
\newcommand{\gl}[1]{Gl.~(\ref{eq:#1})}
\newcommand{\vgl}[1]{vgl. Abschnitt \ref{sec:#1}}

% Anmerkungen am Seitenrand während des Schreibens (für fehlende Literatur etc.)
\newcommand{\anmerkung}[1]{\marginpar{\framebox[14mm][c]{\tiny #1}}}
\newcommand{\suchen}{\anmerkung{Lit. suchen!}}
\newcommand{\kommentar}[1]{\texttt{[#1]}}

% Häufige benötigte Formelelemente
\newcommand{\inv}{^{-1}}
\newcommand{\Matrix}[2]{\left[\!\!\begin{array}{#1} #2 \end{array}\!\!\right]}
\newcommand{\Vektor}[1]{\Matrix{c}{#1}}

\newcommand{\bs}{\boldsymbol}
\newcommand{\ul}{\underline}
\newcommand{\mgl}{$\,=\,$}
\newcommand{\mapp}{$\,$\approx$\,$}
\newcommand{\einh}[1]{\,\mathrm{#1}}
\newcommand{\komma}{,\!}

\newcommand\declare[2]{
 \begin{tabular}{p{15mm}l}
  $#1$ & #2\\
 \end{tabular}
}
\newcommand\define[3]{
 \newcommand{#1}{#2}
 \declare{#2}{#3}
}
\newcommand\redefine[3]{
 \renewcommand{#1}{#2}
 \declare{#2}{#3}
}

\textit{Selten benutzte Formelzeichen und Begriffe sowie abweichende
Bedeutungen werden ausschließlich in Text erläutert.}

\subsection*{Lateinische Notation}
 \define   {\D}   {D}              {Dämpfungsgrad}\\
 \define   {\M}   {\ul{\bs{M}}}    {Massenmatrix}\\
 \define   {\q}   {\bs{q}}         {generalisierte Auslenkung}\\
 \redefine {\dq}  {\bs{\dot q}}    {generalisierte Geschwindigkeit}\\
 \define   {\ddq} {\bs{\ddot q}}   {generalisierte Beschleunigung}\\
 \define   {\Q}   {\bs{Q}}

\subsection*{Griechische Notation}
 \declare   {\alpha}   {Ein Winkel} \\
 \declare   {\Phi}     {Was Anderes}
